\documentclass[10pt,a4paper]{amsart}
\usepackage[margin=19mm]{geometry}
\usepackage{amsmath,amssymb,mathtools}
\usepackage[scr=rsfs,cal=euler]{mathalfa}
\usepackage{xfrac}
\usepackage[unicode,hidelinks,bookmarks=false]{hyperref}
\newcommand{\ie}{i.e.\ }
\newcommand{\cf}{cf.\ }
\newcommand{\resp}{respectively\ }
\newcommand{\Subsection}{\subsection}
\providecommand{\nobreakdash}{\nobreak\hbox{-}}
\setlength{\emergencystretch}{2em}
\multlinegap0pt
\usepackage{bm}
\usepackage{bbm}
\usepackage{dsfont}
\usepackage{xspace}
\usepackage{cancel}
\usepackage{mathtools}
\usepackage{amsthm}
\makeatletter
\@ifundefined{equation}{\newtheorem{equation}{Equation}}{}
\@ifundefined{dcases*}{\newtheorem{dcases*}{Dcases*}}{}
\@ifundefined{lemma}{\newtheorem{lemma}[equation]{Lemma}}{}
\makeatother
\DeclareMathOperator{\Lip}{Lip}
\DeclareMathOperator{\SL}{SL}
\DeclareMathOperator{\arsinh}{arsinh}
\newcommand{\C}{\mathbb{C}}
\newcommand{\Hb}{\mathbb{H}}
\newcommand{\R}{\mathbb{R}}
\newcommand{\dee}{\partial}
\newcommand{\e}{\varepsilon}
\title[Quantitative Equidistribution on Hyperbolic Surfaces and Ari]{Quantitative Equidistribution on Hyperbolic Surfaces and Arithmetic Applications}
\author{Peter Humphries}
\date{Source: arXiv:2512.15664v2 (CC BY 4.0)}
\begin{document}
\maketitle

\noindent\emph{Source and licence.} Quantitative Equidistribution on Hyperbolic Surfaces and Arithmetic Applications. Version arXiv:2512.15664v2 (CC BY 4.0), distributed under the Creative Commons Attribution 4.0 International licence, as recorded in the arXiv licence metadata for this exact version and shown on the abstract page https://arxiv.org/abs/2512.15664v2. Full rights notice: licenses/arxiv-2512.15664-CC-BY-4.0.txt.\par\smallskip

\noindent\textit{Editorial scope.} Counted unit: the Lemma whose source label is \texttt{lem:smoothing} in the article (source0.tex lines 396--406, source label \texttt{lem:smoothing}), delivered together with the article's own complete original proof of it (source0.tex lines 408--458, one \texttt{proof} environment of 51 lines). No mathematics is written, repaired or re-derived: statement and proof are the article's own text line for line. Required context retained uncounted: eqn:geodesicpolar (lines 385--388). Every definition, display and label of those lines is retained unchanged. Omitted from this package and not redistributed: 1--384, 389--395, 407--407, 459--761. Nothing inside a retained excerpt is omitted or altered: every retained body is delivered exactly as the article's lines are written, so no omission receipt of any kind accompanies this package. Citations: the retained excerpts carry no citation command at all, so no bibliography entry of the article is transcribed and no reference is redistributed. Reference adaptation: none was needed and none was made; every reference inside the delivered unit resolves inside it, so no unresolved reference or citation is produced. Printed slips kept literal: no printed word, symbol, hypothesis or number of the counted statement or of its proof was altered. Layout and class: the article's own class is \texttt{amsart} with packages including geometry, amsmath, amsfonts, amsthm, amssymb, enumitem, mathrsfs, mathtools, microtype, hyperref, xcolor; the wrapper is the lane's pinned \texttt{amsart} editorial wrapper at 10pt on A4, which restates the theorem-like environments the retained text needs and adds the article's own macro definitions that the retained text needs (\texttt{\textbackslash C}, \texttt{\textbackslash Hb}, \texttt{\textbackslash Lip}, \texttt{\textbackslash R}, \texttt{\textbackslash SL}, \texttt{\textbackslash arsinh}, \texttt{\textbackslash dee}, \texttt{\textbackslash e}), with extra wrapper packages (bm, bbm, dsfont, xspace, cancel, mathtools). The delivered numbering is the unit's own. Assets: no retained span contains a figure environment, a graphics-inclusion command, a PDF-page inclusion, a table or a TikZ picture; no image file is copied into this package, the package declares no asset dependency and no image file is redistributed with it. Licence: version arXiv:2512.15664v2 of this article is distributed under the Creative Commons Attribution 4.0 International licence, as recorded in the arXiv licence metadata for this exact version and shown on the abstract page https://arxiv.org/abs/2512.15664v2. A full rights notice is delivered at licenses/arxiv-2512.15664-CC-BY-4.0.txt. Characters: every retained excerpt is pure ASCII, so the delivered engine, XeLaTeX with its default Latin Modern text font, reports no missing character.
\begin{equation}
\label{eqn:geodesicpolar}
x = \frac{2 \sqrt{u(u + 1)} \sin \theta}{1 + 2u + 2 \sqrt{u(u + 1)} \cos \theta}, \qquad y = \frac{1}{1 + 2u + 2 \sqrt{u(u + 1)} \cos \theta},
\end{equation}

\begin{lemma}
\label{lem:smoothing}
Let $\Gamma$ be a lattice in $\SL_2(\R)$. Given $F \in \Lip_1(\Gamma \backslash \Hb)$ and $\e > 0$, there exists some $F_{\e} \in C^{\infty}(\Gamma \backslash \Hb)$ for which
\begin{align}
\label{eqn:Fesup}
\sup_{z \in \Gamma \backslash \Hb} |F(z) - F_{\e}(z)| & \leq \e,	\\
\label{eqn:nablaFesup}
\sup_{z \in \Gamma \backslash \Hb} \Im(z)^2 \left|\frac{\dee F_{\e}}{\dee z}\right|^2 & \leq \left(e^{\e} - \frac{1}{2}\right)^2,
\end{align}
where $\frac{\dee}{\dee z} \coloneqq \frac{1}{2} (\frac{\dee}{\dee x} - i \frac{\dee}{\dee y})$ denotes the Wirtinger derivative.
\end{lemma}

\begin{proof}
For fixed $\e > 0$, let $k_{\e} : \R_+ \to \C$ be the smooth nonnegative function
\[k_{\e}(u) \coloneqq \begin{dcases*}
\frac{1}{C\sinh^2 \frac{\e}{2}} \exp\left(\frac{\sinh^2 \frac{\e}{2}}{u^2 - \sinh^2 \frac{\e}{2}}\right) & if $0 < u < \sinh^2 \frac{\e}{2}$,	\\
0 & if $u \geq \sinh^2 \frac{\e}{2}$,
\end{dcases*}\]
where $C \coloneqq 4\pi \int_{0}^{1} \exp(\frac{1}{u^2 - 1}) \, du$. By passing to geodesic polar coordinates, we have that for all $z \in \Hb$,
\begin{equation}
\label{eqn:keint}
\int_{\Hb} k_{\e}(u(z,w)) \, d\mu(w) = 4\pi \int_{0}^{\infty} k_{\e}(u) \, du = 1.
\end{equation}
We let $K_{\e} : \Gamma \backslash \Hb \times \Gamma \backslash \Hb \to \C$ denote the automorphic kernel
\[K_{\e}(z,w) \coloneqq \sum_{\gamma \in \Gamma} k_{\e}(u(z,\gamma w)).\]

We now define the desired function $F_{\e} : \Hb \to \C$ via
\[F_{\e}(z) \coloneqq \int_{\Gamma \backslash \Hb} F(w) K_{\e}(z,w) \, d\mu(w) = \int_{\Hb} F(w) k_{\e}(u(z,w)) \, d\mu(w).\]
The smoothness of $F_{\e}$ is clear from the smoothness of $k_{\e}$, while the $\Gamma$-invariance of $F_{\e}$ follows from the bi-$\Gamma$-invariance of $K_{\e}$.

Next, for any $z \in \Hb$, we have via \eqref{eqn:keint} that
\begin{align*}
\left|F(z) - F_{\e}(z)\right| & = \left|\int_{\Hb} (F(z) - F(w)) k_{\e}(u(z,w)) \, d\mu(w)\right|	\\
& \leq \int_{\Hb} \rho(z,w) k_{\e}(u(z,w)) \, d\mu(w)	\\
& = \int_{\Hb} \rho(i,w) k_{\e}(u(i,w)) \, d\mu(w)
\end{align*}
as $F$ is $1$-Lipschitz and $u,\rho,\mu$ are $\SL_2(\R)$-invariant. Passing to geodesic polar coordinates, observing that $\rho(i,w) = 2\arsinh\sqrt{u(i,w)} \leq \e$ for $u(i,w) \leq \sinh^2 \frac{\e}{2}$, and recalling \eqref{eqn:keint}, we thereby obtain \eqref{eqn:Fesup}.

When $\frac{\dee F_{\e}}{\dee z}$ is nonzero, its modulus is equal to half of the modulus of $\nabla_{v_0} F_{\e}(z)$, the directional derivative of $F_{\e}$ in the direction
\[v_0 = \frac{\frac{\dee F_{\e}}{\dee z}}{\left|\frac{\dee F_{\e}}{\dee z}\right|}.\]
Thus if $\frac{\dee F_{\e}}{\dee z} \neq 0$, then
\[\Im(z)^2 \left|\frac{\dee F_{\e}}{\dee z}\right|^2 = \frac{1}{4} \Im(z)^2 \left|\nabla_{v_0} F_{\e}(z)\right|^2 = \frac{1}{4} \Im(z)^2 \lim_{h \to 0} \left|\frac{F_{\e}(z + hv_0) - F_{\e}(z)}{h}\right|^2.\]
Let $g_0,g_1 \in \SL_2(\R)$ be such that $g_0 i = z$ and $g_1 i = z + hv_0$. Then by making the change of variables $w \mapsto g_0 w$ and $w \mapsto g_1 w$ respectively together with the facts that $u$ and $\mu$ are $\SL_2(\R)$-invariant and $F$ is $1$-Lipschitz, we see that
\begin{align*}
\left|F_{\e}(z + hv_0) - F_{\e}(z)\right| & = \left|\int_{\Hb} F(w) \left(k_{\e}(u(g_1 i,w)) - k_{\e}(u(g_0 i,w))\right) \, d\mu(w)\right|	\\
& = \left|\int_{\Hb} \left(F(g_1 w) - F(g_0 w)\right) k_{\e}(u(i,w)) \, d\mu(w)\right|	\\
& \leq \int_{\Hb} \rho(g_1 w,g_0 w) k_{\e}(u(i,w)) \, d\mu(w).
\end{align*}
Since $g_1 i = z + hv_0$ and $g_0 i = z$, we must have that $g_1 w = \Re(z + hv_0) + \Im(z + hv_0)w$ and $g_0 w = \Re(z) + \Im(z)w$. Thus
\[\rho(g_1 w,g_0 w) = 2 \arsinh \frac{|h| \left|v_0 + \Im(v_0) (w - i)\right|}{2\sqrt{\Im(z)(\Im(z) + h \Im(v_0))} \Im(w)},\]
so that
\[\lim_{h \to 0} \frac{\rho(g_1 w,g_0 w)}{|h|} = \frac{\left|v_0 + \Im(v_0) (w - i)\right|}{\Im(z) \Im(w)} \leq \frac{1 + |w - i|}{\Im(z) \Im(w)}.\]
It follows that
\[\Im(z)^2 \left|\frac{\dee F_{\e}}{\dee z}\right|^2 \leq \frac{1}{4} \left(\int_{\Hb} \frac{1 + |w - i|}{\Im(w)} k_{\e}(u(i,w)) \, d\mu(w)\right)^2.\]
Recall that $k_{\e}(u(i,w))$ vanishes unless $u(i,w) \leq \sinh^2 \frac{\e}{2}$. When this inequality holds, we see via passing once more to geodesic polar coordinates as in \eqref{eqn:geodesicpolar} that
\begin{align*}
\frac{1 + |w - i|}{\Im(w)} & = 1 + 2u + 2 \sqrt{u(u + 1)} \cos \theta + 2\sqrt{u \left(1 + 2u + 2\sqrt{u(u + 1)} \cos \theta\right)}	\\
& \leq 1 + 2 \sinh^2 \frac{\e}{2} + 2 \sinh \frac{\e}{2} \cosh \frac{\e}{2} + 2 \sinh \frac{\e}{2} \sqrt{1 + 2 \sinh^2 \frac{\e}{2} + 2 \sinh \frac{\e}{2} \cosh \frac{\e}{2}}	\\
& = \cosh \e + \sinh \e + 2 \sinh \frac{\e}{2} \sqrt{\cosh \e + \sinh \e}	\\
& = 2e^{\e} - 1.
\end{align*}
With this inequality in hand, the bound \eqref{eqn:nablaFesup} then follows once more from \eqref{eqn:keint}.
\end{proof}

\end{document}
